\section{Reproducibility Details}
\label{app:reproducibility}

\paragraph{Fine-tuning configuration.}
All four fine-tuned models (Qwen3-8B and Llama-3.1-8B-Instruct on SciERC and DocRED) use identical LoRA hyperparameters: rank $r{=}16$, $\alpha{=}32$, dropout $0.05$, no bias, targeting all attention and MLP projection layers (\texttt{q\_proj}, \texttt{k\_proj}, \texttt{v\_proj}, \texttt{o\_proj}, \texttt{gate\_proj}, \texttt{up\_proj}, \texttt{down\_proj}).
Training uses AdamW with learning rate $2{\times}10^{-4}$, cosine schedule, warmup ratio $0.1$, effective batch size $16$, bf16 precision, gradient checkpointing, and completion-only loss for 3 epochs.
SciERC trains on 1,762 examples (260 dev) for 333 steps; DocRED trains on 3,053 examples (998 dev) for 573 steps.
Maximum sequence length is 2,048 (SciERC) and 4,096 (DocRED).
All runs use seed 42.

\paragraph{Hardware and software.}
Training runs on a single NVIDIA RTX PRO 6000 Blackwell Server Edition GPU (98\,GB video RAM (VRAM)).
SciERC fine-tuning takes approximately 11 minutes per model; DocRED takes approximately 1.5 hours.
Software versions: PyTorch 2.10, Transformers 4.57, PEFT 0.19, TRL 1.1, Accelerate 1.13.

\paragraph{Inference time.}
Inference timing is reported in Appendix~\ref{app:runtime} (NVIDIA RTX PRO 6000 Blackwell Server Edition GPU, same hardware as training).
S1 constrained decoding adds $+$13.1\% wall-clock overhead on SciERC (Qwen3-8B, zero-shot greedy), lower on DocRED due to amortization over longer sequences; under LoRA fine-tuning, overhead is $+$10.0\% on SciERC.

\paragraph{Greedy decoding determinism.}
\label{app:greedy-determinism}
All main results use greedy decoding (temperature$\,{=}\,$0, argmax token selection).
Argmax is deterministic: the random seed controls \texttt{torch.manual\_seed} for sampling and dropout, neither of which is active during greedy inference.
We verify this with three seeds (42, 123, 456) across two LoRA fine-tuned models on all 551 SciERC documents.
Seeds 123 and 456, evaluated on the same GPU, produce byte-identical outputs for both Qwen3-8B and Llama-3.1-8B.
The residual variance in Table~\ref{tab:seed-variance} stems entirely from cross-hardware floating-point divergence: seed 42 was evaluated on a different server with a different CUDA driver version.
Even under this conservative upper bound, $\sigma(\text{E-F1}) \leq 0.16$\,pp, $\sigma(\text{R-F1}) \leq 0.43$\,pp, and $\sigma(\text{DR\%}) \leq 0.08$\,pp.
The constrained decoding DR\% reduction ($-$2.06\,pp for Qwen, $-$1.01\,pp for Llama) exceeds this variance by ${>}12\times$ (signal-to-noise ratio ${>}$\,12), confirming that observed improvements reflect genuine method effects rather than run-to-run noise.
We therefore report single-run results without error bars; bootstrap confidence intervals (10{,}000 iterations) are provided for F1 metrics.

\begin{table}[h]
\centering
\caption{Cross-seed variance for LoRA fine-tuned models on SciERC (551 documents, greedy decoding, seeds 42/123/456). Values: mean $\pm$ population $\sigma$ across three seeds. Seeds 123 and 456 produce identical outputs on the same GPU; all observed $\sigma$ reflects cross-server floating-point divergence.}
\label{tab:seed-variance}
\small
\resizebox{\columnwidth}{!}{%
\begin{tabular}{@{}lllccc@{}}
\toprule
Model & Dataset & Mode & E-F1 & R-F1 & DR\% \\
\midrule
Qwen3-8B & SciERC & Unc & $66.94 \pm 0.08$ & $43.25 \pm 0.43$ & $2.06 \pm 0.08$ \\
Qwen3-8B & SciERC & Con & $66.94 \pm 0.08$ & $43.22 \pm 0.39$ & $0.00 \pm 0.00$ \\
Llama-3.1-8B & SciERC & Unc & $67.99 \pm 0.16$ & $45.38 \pm 0.05$ & $1.04 \pm 0.05$ \\
Llama-3.1-8B & SciERC & Con & $67.99 \pm 0.16$ & $45.40 \pm 0.05$ & $0.03 \pm 0.04$ \\
\bottomrule
\end{tabular}}
\end{table}

\paragraph{Beam search non-determinism.}
Greedy decoding is deterministic (\S\ref{app:greedy-determinism} above).
Beam search ($k{\geq}2$) introduces non-determinism from GPU floating-point accumulation order; we report single-run results.
Based on the greedy cross-hardware variance ($\sigma < 0.5$\,pp, Table~\ref{tab:seed-variance}), beam search variance is expected to be of similar magnitude.

\paragraph{200-document subset justification.}
\label{app:power-analysis}
The 70B LoRA evaluation on DocRED uses a 200-document subset (vs.\ the full 998-document dev set used for 8B models).
To verify that this subset provides sufficient statistical power, we conduct a bootstrap analysis using the full 998-document 8B LoRA results as a reference population.
Resampling 200 documents 5{,}000 times yields a relation F1 margin of error of $\pm$3.3 points (95\% CI) and a DR violation rate margin of error of $\pm$0.45 percentage points---precise enough to support our conclusions.
For zero-shot models, where DR violations are prevalent (15.0\% unconstrained), a 200-document subset achieves 100\% power to detect the constrained decoding effect (Cohen's $h = 0.80$); for fine-tuned models, where violations are already rare ($<$0.5\%), 94.9\% of 200-document subsamples correctly identify the direction of DR reduction, and the primary effect is established on the full 998-document evaluation.
